\scanpage{55}
\begin{center}- 130 -\end{center}
\begin{flushright}IX\end{flushright}

Elle implique bien que si $\varphi^t_{\xi_{\unsure{o|0}}}$ est symétrique et positive, il en est de même de $\varphi^t_\xi$.

Enfin, le fait que $\varphi^t_{\xi_{\unsure{o|0}}}$ soit non dégénérée équivaut évidemment à l'assertion que l'accouplement (10.4.1) est séparant (compte tenu que
$\xi_{\unsure{ot|0t}}:M\otimes_{\mathbb Z}\mathbb Q\longrightarrow M'\otimes_{\mathbb Z}\mathbb Q$
est un isomorphisme, $\widetilde\xi_{\unsure{o|0}}$ étant une isogénie), et est donc également indépendant de la polarisation choisie. On peut d'ailleurs noter que pour vérifier que (10.4.1) est séparant, il suffit de vérifier la même assertion pour l'accouplement de monodromie $u_\ell$ (9.1.2) pour \uline{un} nombre premier $\ell$. Prenant $\ell\ne p$, on peut utiliser la description 9.2 de $u_\ell$; le fait que $u_\ell$ soit séparant à gauche (donc des deux côtés, pour une raison de rangs) est alors trivial.

11. \uline{Composantes connexes du modèle de Néron}: \uline{relation de dualité, comportement asymptotique}.

11.0. Le but du présent paragraphe est l'étude du $k$-groupe fini étale
\[
\text{(11.0.1)}\qquad
\Phi_{\unsure{o|0}}=A_{\unsure{o|0}}/A_{\unsure{o|0}}^{\unsure{o|0}}
\]
déjà considéré dans 1. Pour tout nombre premier $\ell$, nous désignons par
\[
\Phi_{\unsure{o|0}}(\ell)
\]
la composante $\ell$-primaire de ce dernier, de sorte que la connaissance de
$\Phi_{\unsure{o|0}}$ équivaut à celle des
$\Phi_{\unsure{o|0}}(\unsure{l|1})$, grâce à la formule
\[
\text{(11.0.2)}\qquad
\Phi_{\unsure{o|0}}=\underset{\ell}{\unsure{\coprod|\prod}}\Phi_{\unsure{o|0}}(\ell).
\]
\margin{\unsure{448|148}}
\margin{448}
\margin{54}
